\paragraph{名目GDP水準目標（NGDPLT）}
\label{sec:chapter5_beta_shock_policies_ngdplt}
式
\eqref{form:chapter3_policies_monetary_home_ngdplt_i_H_eq}、
\eqref{form:chapter3_policies_monetary_home_ngdplt_gamma_H_eq}および
\eqref{form:chapter3_policies_monetary_home_ngdplt_chi_H_eq}
により記述されたとおり、NGDPLTの式は以下のようであった。
\begin{align}
\label{form:chapter5_beta_shock_policies_ngdplt_i_H_eq}
    i_t^H&=i_{ss}^H+\phi_{gap}^H(\log(p_t^H y_t^H)-\log(\chi_t^H))+\phi_{level}^H\gamma_t^H \\
\label{form:chapter5_beta_shock_policies_ngdplt_gamma_H_eq}
    \gamma_t^H&=\gamma_{t-1}^H+(\log(p_{t-1}^H y_{t-1}^H)-\log(\chi_{t-1}^H)) \\
\label{form:chapter5_beta_shock_policies_ngdplt_chi_H_eq}
    \log(\chi_t^H)-\log(\chi_{ss}^H)&=\rho_{\chi}^H\left(\log(\chi_{t-1}^H)-\log(\chi_{ss}^H)\right)+\varepsilon_t^{\chi,H}
\end{align}
ここで
\(\phi_{gap}^H, \phi_{level}^H\)は反応係数、
\(p_t^H y_t^H\)は名目GDP水準、
\(\chi_t^H\)は目標パス、
\(\gamma_t^H\)は過去の乖離の累積、
\(\varepsilon_t^{\chi,H}\)は外生ショック項である。
\begin{enumerate}
    \item
    \label{itm:chapter5_beta_shock_policies_ngdplt_c_H_W}
        \(c_t^{H\to W}\)：
        総消費指数\(c_t^{H\to W}\)の図
        \ref{fig:chapter5_beta_shock_graphs_c_H_W}
        を見ると、NGDPLTの軌道は生産水準目標（OLT）とほぼ完全に重なっていることがわかる。
        100期までの前半において、NGDPLTの落ち込みは
        総消費水準目標（CLT）や名目総消費水準目標（NCLT）には及ばないものの、
        消費者物価水準目標（CPLT）や消費者物価インフレ目標（IT-CPI）のような極端な大暴落を免れている。
        一方で中盤以降のプラス圏での推移を見ると、厚生を改善させる消費の増加幅は小規模にとどまっている。
        CPLTやIT-CPIが描く巨大な山や、NCLTが長期にわたって維持する高い水準と比較すると、
        消費の増加による厚生の押し上げ効果は限定的となっている。
    \item
    \label{itm:chapter5_beta_shock_policies_ngdplt_y_H}
        \(y_t^H\)：
        生産\(y_t^H\)の図
        \ref{fig:chapter5_beta_shock_graphs_y_H}
        を見ると、NGDPLTの軌道は総消費指数\(c_t^{H\to W}\)と同様に
        生産水準目標（OLT）とほぼ重なっていることがわかる。
        100期までの前半において、NGDPLTの落ち込みは
        総消費水準目標（CLT）や名目総消費水準目標（NCLT）よりは大きいが、
        消費者物価水準目標（CPLT）や消費者物価インフレ目標（IT-CPI）のような極端な落ち込みには至っていない。
        生産の低下は労働の減少を意味するため厚生にはプラスに働くが、
        NGDPLTにおける前半の厚生押し上げ効果はCPLTやIT-CPIほど大きくはない。
        一方で中盤以降のプラス圏での推移（オーバーシュート）を見ると、
        生産の増加は労働の非効用を通じて厚生を悪化させる要因となる。
        NGDPLTのオーバーシュートはCLTやNCLTよりも低く抑えられており、
        またCPLTやIT-CPIのような遅くて巨大な山を作ることもないため、
        後半の生産過剰に伴う厚生のマイナス効果は比較的小さくとどまっている。
    \item
    \label{itm:chapter5_beta_shock_policies_ngdplt_pi_H}
        \(\pi_t^H\)：
        グロス・インフレ率\(\pi_t^H\)の図
        \ref{fig:chapter5_beta_shock_graphs_pi_H}
        を見ると、NGDPLTの定常状態からの乖離は、
        生産者物価水準目標（PPLT）や生産者物価インフレ目標（IT-PPI）、名目総消費水準目標（NCLT）に次いで
        極めて小さく抑えられていることがわかる。
        消費者物価インフレ目標（IT-CPI）のような極端な初期の下落や、
        総消費水準目標（CLT）のような大きな変動は見られず、
        全期間を通じてゼロ近傍で極めて安定した推移を維持している。
        名目GDPという名目アンカーを持つことでインフレ率のブレがしっかりと制御されており、
        この定常状態からの乖離の小ささがフローの厚生損失を最小限にとどめることに貢献している。
    \item
    \label{itm:chapter5_beta_shock_policies_ngdplt_Delta_t_minus_1_H}
        \(\Delta_{t-1}^H\)：
        価格分散\(\Delta_{t-1}^H\)の図
        \ref{fig:chapter5_beta_shock_welfare_Delta_H}
        を見ると、NGDPLTの定常状態からの乖離は、
        生産者物価水準目標（PPLT）や生産者物価インフレ目標（IT-PPI）、名目総消費水準目標（NCLT）と並んで
        全政策の中で最も小さく抑えられていることがわかる。
        総消費水準目標（CLT）のような著しい増大や、
        消費者物価水準目標（CPLT）および消費者物価インフレ目標（IT-CPI）のような大きな山を作ることはなく、
        全期間を通じてほぼゼロの水準を維持している。
        これは前述のインフレ率の極めて安定した推移が価格分散の蓄積を未然に防いでいる結果であり、
        この生産効率の悪化の回避がNGDPLTの厚生を高く保つ重要な要因となっている。
\end{enumerate}
以上の比較分析から、NGDPLTはNCLTほどの需要回復力は持たないものの、
実体経済の過熱を抑制し、最も安定的かつ早期に経済を平時へと帰還させる能力に長けていると言える。
この性質が、厚生$W^H$評価においてNCLTに次ぐ高い評価を支える要因となっている。